% SPDX-FileCopyrightText: 2026 Frederic Ruget <fred@atlant.is> (GitHub: @douzebis)
%
% SPDX-License-Identifier: CC-BY-4.0
\documentclass[11pt,a4paper]{article}
\usepackage[T1]{fontenc}
\usepackage[utf8]{inputenc}
\usepackage{lmodern}
\usepackage[margin=28mm]{geometry}
\usepackage{amsmath,amssymb,amsthm}
\usepackage{booktabs}
\usepackage{microtype}
\usepackage[hidelinks]{hyperref}

\newtheorem{theorem}{Theorem}
\newtheorem{proposition}[theorem]{Proposition}
\newtheorem{corollary}[theorem]{Corollary}
\newtheorem{lemma}[theorem]{Lemma}
\theoremstyle{definition}
\newtheorem{definition}[theorem]{Definition}
\newtheorem{conjecture}[theorem]{Conjecture}
\newtheorem{question}[theorem]{Question}
\theoremstyle{remark}
\newtheorem{remark}[theorem]{Remark}

\newcommand{\lean}[1]{\texttt{#1}}
\newcommand{\str}[1]{\mathtt{#1}}
\newcommand{\C}{\mathbf{C}}
\newcommand{\repo}{\url{https://github.com/douzebis/blog.atlant.is}}

\title{Substitutions and resets in Levine's hat game:\\
a structure theorem and new lower bounds\\
for two, three and four players}
\author{Fr\'ed\'eric Ruget}
\date{November 2026 (draft)}

\begin{document}
\maketitle

\begin{abstract}
In Levine's hat game each of $n$ players wears a stack of hats, black or
white at random; each sees the stacks of the others and points at one hat
of their own stack; they win when all the hats pointed at are black. We
give an exact formula for the number of winning positions of a strategy
obtained by putting one strategy in the place of a hat of another
(\emph{substitution}), of which the resets and shifts used so far are
cases. For two players it yields in one line a known family of strategies
with $\lfloor 7\cdot 4^h/20\rfloor$ winning pairs on $h$ hats, a second way to add two hats with the best possible
gain, and a reduction of the conjecture $V_2 = 7/20$ to finite stacks.
With the same tools we improve the best known lower bounds: for three
players, $V_3 \ge 73225/261950 \approx 0.279538$; for four players,
$V_4 \ge 15191/64120 \approx 0.236915$; and for two players with hats
black with probability $p \ne 1/2$, a bound that beats the published
ones for most $p$. The statements, and the counts of the tables for three
and four players, are checked in the Lean proof assistant.
\end{abstract}

\noindent\textbf{MSC 2020:} 91A12 (primary); 05A05, 60C05, 68V20
(secondary).

\noindent\textbf{Keywords:} hat puzzles, Levine's hat problem,
cooperative games, lower bounds, formal proof, Lean.

\section{Introduction}

\subsection{The game}

Two players, Alice and Bob, each wear a stack of hats. Each hat is black
or white, by the toss of a fair coin. Alice sees Bob's stack and not her
own, and Bob sees Alice's. Without a word, each points at one hat of their
own stack. They win when both hats are black. They may agree on a plan
beforehand. The game is due to Levine (2011); it became known through
\cite{Kh11}, and \cite{BFG} is the reference on it. Write $V_2$ for the
best probability of winning, over all plans and all heights of the
stacks; and $V_n$ for the game with $n$ players, who must all point at a
black hat, each seeing the stacks of all the others.

Pointing at random wins one time in four. The first surprise is that one
can do better; the second, that nobody knows how much better. The known
bounds for two players are
\[
  \frac{7}{20} \;\le\; V_2 \;\le\; \frac{81}{224} ,
\]
both from \cite{BFG}, and $V_2 = 7/20$ is conjectured. The upper bound
was lowered to $57/160$ and, according to a recent computer-assisted
note that has not been reviewed, to $17/48$ \cite{Ya26}; an analytic
bound without computer is in \cite{HT25}. For more players, $V_n$
decreases strictly with $n$ \cite{FKK21,AFKK}, and \cite{BFG} gives
$V_3 \ge 9119/32670$ and $V_4 \ge 14844/64120$, from tables found by a
search, due to Reyes and Carter.

\subsection{Earlier constructions}

The strategies that reach $7/20$ have a short history, part of which is
not in journals.

In 2014 the notes of Wagon \cite{Wa14} reported tables found by Carter
and Reyes on up to eight hats, with $1, 5, 22, 89, 358, 1433, \ldots$
winning pairs, and the idea of a \emph{reset}, due to Carter, Reyes and
Rosenberg: when the lowest hats seen make the game hopeless or settled,
start again on the hats above. Resetting a table on three hats with $22$
winning pairs gives exactly $7/20$ (Tiefenbruck, Reyes, Carter and
Rosenberg); this is the lower bound of \cite{BFG}. The book \cite{VW20}
states, with a sketch of proof, that there is a strategy with
$\lfloor 7\cdot 4^h/20\rfloor$ winning pairs on every number $h$ of hats.

In 2020 the author described these strategies on a blog, in three posts
\cite{R20}, without knowing the book. The posts introduce operations
that act on arbitrary strategies, each with an exact formula for the
number of winning pairs of the result: three \emph{resets} (on all
white, on all black, on both), and three \emph{shifts}, which add one
hat or two and change the choice for one or two stacks. They name
\emph{Carter strategies} the family $\C_h$ obtained from the strategy on
one hat by shifts alone, because it gives on few hats the scores of
Carter's tables, and \emph{Reyes strategies} the resets of $\C_3$ by
itself. These names and operations are used below; they are defined
again in Section~\ref{sec:two}, so that the paper can be read alone.

In 2025, \cite{BCCZW} gave a ``combining construction'', which is the
reset on both colors, and strategies with the scores of $\C_h$ for odd
$h$. The present paper replaces the separate formulas of the posts of
2020 by one theorem, of which each is a case; a longer account, which
revises those posts, is \cite{web}.

\subsection{Results}

A \emph{strategy} on $h$ hats is a table that gives, for each stack seen,
the position to point at. Good strategies are built from small ones by a
few operations: \cite{BFG} uses \emph{resets}, which start the game again
on the hats above when the lowest hats are all white or all black, and
\cite{BCCZW} a ``combining construction''. Our first result is one
formula that holds all of these.

\begin{itemize}
\item \textbf{Substitution} (Theorem~\ref{thm:sub}). Put a strategy $s$
  on $h$ hats in the place of the bottom hat of a strategy $m$ on $k+1$
  hats. The result has $4^k(|H(s)|-1)+|H(m)|$ winning pairs of stacks,
  where $|H(\cdot)|$ counts the winning pairs. So every strategy $m$ is
  an \emph{operation} that adds $k$ hats, with a gain that does not
  depend on $s$. Resets and shifts are substitutions
  (Proposition~\ref{prop:instances}).
\item \textbf{Consequences for two players} (Section~\ref{sec:two}).
  The strategies $\C_h$ with
  $\lfloor 7\cdot 4^h/20\rfloor$ winning pairs on $h$ hats, for every
  $h$, which are stated with a sketch of proof in \cite{VW20}, follow in
  one line, with a complete proof; so does the remark of Yeroshkin
  \cite{VW20} that the conjecture $V_2 = 7/20$ holds if and only if no
  strategy on any number of hats wins one more pair than $\C_h$; an operation that
  adds $k$ hats gains at most the best score on $k+1$ hats less $4^k$;
  and there is a second operation that adds two hats with the best gain,
  not equivalent to the first.
\item \textbf{Biased coins} (Section~\ref{sec:bias}). When hats are
  black with probability $p$, a strategy that plays the best strategy on
  three hats on the first block of three hats that is not all white wins
  with probability
  $(3p^2-6p^3+8p^4-6p^5+2p^6)/(1-(1-p)^6)$. For $p$ from $1/20$ to
  $8/20$ this is more than the strategies of \cite{BFG} and
  \cite{BCCZW}; and $V_2(1-p) = V_2(p) + 1 - 2p$ carries the result to
  $p > 1/2$.
\item \textbf{More players} (Section~\ref{sec:players}). The reset of
  \cite{BFG} for $n$ players has an exact count of winning positions
  (Theorem~\ref{thm:reset}), from which the limit formula of \cite{BFG}
  follows. A search that starts from a table made by a reset, instead of
  a random table, then gives
  \[
    V_3 \;\ge\; \frac{73225}{261950} = 0.279538\ldots ,
    \qquad
    V_4 \;\ge\; \frac{15191}{64120} = 0.236915\ldots
  \]
  against $0.279124\ldots$ and $0.231503\ldots$ in \cite{BFG}.
\end{itemize}

Section~\ref{sec:upper} relates the upper bounds on $V_2$ obtained from
``hint games'' to the strategies $\C_h$, and
Section~\ref{sec:programs} reports what exhaustive searches say of the
best strategies on up to five hats.

\subsection{What is checked, and how}

Every numbered theorem, proposition and corollary of
Sections~\ref{sec:game} to~\ref{sec:players} is stated and proved in the
Lean~4 proof assistant with Mathlib; the names of the Lean theorems are
given after each statement and collected in
Section~\ref{sec:lean}. This includes the two bounds displayed above:
the tables are part of the Lean sources, and the kernel of Lean counts
their winning positions one by one, so that no program has to be
trusted. Where a statement rests on a program instead, the text says so.

A longer account, with the history of these strategies since 2020,
pictures and a calculator, is on the web \cite{web}.

\section{The game, strategies, hits}\label{sec:game}

A \emph{stack} of $h$ hats is a word $t = t_0 t_1 \cdots t_{h-1}$ on
$\{0,1\}$; $t_0$ is the bottom hat, $1$ is black. We also read $t$ as the
number $\sum_i t_i 2^i$. A \emph{strategy} on $h$ hats is a map $s$ from
stacks to positions $\{0,\ldots,h-1\}$: the player who \emph{sees} the
stack $u$ points at position $s(u)$ of their own stack. We write a
strategy as the word $s(0)s(1)\cdots s(2^h-1)$.

Unless said otherwise both players use the same strategy. A pair of
stacks $(t,u)$ is a \emph{hit} of $s$ when $t_{s(u)} = 1$ and
$u_{s(t)} = 1$. We write $H(s)$ for the set of hits; the win rate of $s$
is $|H(s)|/4^h$. A strategy on $h$ hats is also a strategy on more hats,
with the same win rate, so $V_2$ is the supremum of the win rates over
all $h$, and no finite stack attains $7/20$: a win rate on $h$ hats is a
multiple of $4^{-h}$.

Three elementary facts are used throughout.

\begin{proposition}\label{prop:basic}
Let $s$ be a strategy on $h$ hats.
\begin{enumerate}
\item The value of $s$ on the all-white stack does not change $|H(s)|$;
  nor does its value on the all-black stack.
\item Renaming the positions of the hats does not change $|H(s)|$.
\item Exchanging the two colors does not change $|H(s)|$.
\end{enumerate}
\end{proposition}

For the third, the \emph{dual} of $s$ is $\bar s(t) = s(\bar t)$, where
$\bar t$ is $t$ with the colors exchanged; a player points at a black hat
for exactly half of the stacks they may wear, whatever they see, and an
inclusion--exclusion gives $|H(\bar s)| = |H(s)|$. Two strategies are
\emph{equivalent} when one is the other with the positions renamed.
(Lean: \lean{hits\_congr\_white}, \lean{hits\_congr\_black},
\lean{hits\_renamed}, \lean{hits\_dual}.) Throughout, short proofs are given in the text; for
the others the text has a sketch and Appendix~\ref{app:proofs} the
complete proof.

\section{Substitution, and two players}\label{sec:two}

\subsection{The theorem}

\begin{definition}\label{def:sub}
Let $s$ be a strategy on $h \ge 1$ hats and $m$ a strategy on $k+1$ hats.
The \emph{substitution} of $s$ for the bottom hat of $m$ is the strategy
$m[s]$ on $h+k$ hats defined as follows, for a stack whose $h$ lowest
hats are $\ell$ and whose other hats are $x$.
\begin{itemize}
\item If the hats of $\ell$ are not all of the same color: $m[s]$ points
  at $s(\ell)$.
\item If they all have the color $c$: let $p = m(cx)$, the choice of $m$
  for the stack made of one hat of color $c$ under the hats $x$. If
  $p = 0$, $m[s]$ points at $s(\ell)$; otherwise it points at
  $h + p - 1$, the hat that $m$ pointed at.
\end{itemize}
\end{definition}

\begin{theorem}[Substitution]\label{thm:sub}
For every strategy $s$ on $h \ge 1$ hats and $m$ on $k+1$ hats,
\[
  |H(m[s])| \;=\; 4^k \bigl(|H(s)| - 1\bigr) + |H(m)| .
\]
\end{theorem}

\begin{proof}
Sort the pairs of stacks by the $h$ lowest hats of each, which are
\emph{mixed} or of a single color.

\emph{Both mixed.} Both players play $s$: each hit of $s$ between mixed
stacks is counted $4^k$ times, once for each choice of the hats above.

\emph{Both of a single color.} Read each stack as a stack of $k+1$ hats.
A player points at a black hat exactly when $m$ does: if $m$ goes up, it
is the same hat; if $m$ points at the bottom hat, the player points among
$h$ hats that all have the color of that bottom hat. These pairs give
$|H(m)|$ hits.

\emph{All white and mixed.} The player who sees the mixed stack points at
one of the $h$ white hats: no hit.

\emph{All black and mixed.} The player who sees the mixed stack points at
a black hat. The other one either plays $s$, and points at a fixed
position among the lowest hats, or goes up to a fixed position above. In
both cases the hat is black for $2^k(2^{h-1}-1)$ of the $2^k(2^h-2)$
stacks whose lower part is mixed. With the $2^k$ stacks that are all
black below and the two orders, that is $2\cdot 4^k(2^{h-1}-1)$ hits.

The same count for $s$ alone gives
$|H(s)| = M + 2(2^{h-1}-1) + 1$, where $M$ is the number of hits of $s$
between mixed stacks and the last $1$ is the pair of all-black stacks.
Hence $|H(m[s])| = 4^k M + 2\cdot 4^k(2^{h-1}-1) + |H(m)|
= 4^k(|H(s)|-1) + |H(m)|$.
\end{proof}

(Lean: \lean{hits\_sub}.) The fourth case is the heart of the proof:
going up or staying makes no difference to the count, which is why any
$m$ will do. To substitute for another hat of $m$ than the bottom one,
rename the positions first; by Proposition~\ref{prop:basic} nothing new
is obtained up to equivalence.

So each strategy $m$ on $k+1$ hats is an operation that adds $k$ hats,
with the \emph{gain} $|H(m)| - 4^k$: it turns $|H(s)|$ into
$4^k|H(s)|$ plus the gain.

\subsection{Resets and shifts are substitutions}

The following operations are those of \cite{R20} (where black and white
are exchanged, and the white and black operations are called left and
right). Let $s$ be a strategy on $h$ hats and $b$ a strategy on $k$ hats.
\begin{itemize}
\item The \emph{white reset} of $s$ by $b$, on $h+k$ hats: play $s$ on
  the $h$ lowest hats, unless they are all white; then play $b$ on the
  $k$ hats above. The \emph{black reset} is the same with ``all black'',
  and the \emph{double reset} with ``all white or all black''. The
  double reset is the reset of \cite{BFG} for two players and the
  combining construction of \cite{BCCZW}; the white reset is in
  \cite{Wa14}.
\item The \emph{white shift} of $s$, on $h+1$ hats: play $s$ on the $h$
  lowest hats, except that $h$ white hats under a black one send to the
  new hat. The \emph{black shift}: the same, except that $h$ black hats
  under a white one send to the new hat.
\item The \emph{double shift} of $s$, on $h+2$ hats: play $s$ on the $h$
  lowest hats, with two exceptions. Seeing $h$ black hats with a white
  hat just above them, point at the first new hat; seeing $h$ white hats
  with a black hat at the top, point at the top hat.
\end{itemize}
Write $\str{0}$ for the strategy on one hat, and recall that a strategy
is written as the word of its values.

\begin{proposition}\label{prop:instances}
Each reset of $s$ by $b$ is the substitution of $s$ into the same reset
of $\str{0}$ by $b$. The white shift, the black shift and the double
shift of $s$ are $m[s]$ for $m = \str{0010}$, $m = \str{0100}$ and
$m = \str{01002120}$. Consequently, for $b$ on $k$ hats,
\begin{align*}
  |H| &= 4^k|H(s)| + |H(b)| && \text{(white and black resets)},\\
  |H| &= 4^k|H(s)| + 4|H(b)| - 4^k && \text{(double reset)},
\end{align*}
and the shifts have $4|H(s)|+1$ hits (white, black) and
$16|H(s)|+6$ hits (double).
\end{proposition}

(Lean: \lean{lr\_eq\_sub}, \lean{rr\_eq\_sub}, \lean{dr\_eq\_sub},
\lean{ls\_eq\_sub}, \lean{rs\_eq\_sub}, \lean{ds\_eq\_sub}; the counts
are also proved directly: \lean{hits\_lr}, \lean{hits\_rr},
\lean{hits\_dr}, \lean{hits\_ls}, \lean{hits\_rs}, \lean{hits\_ds}.)

\subsection{The Carter strategies}

Following \cite{R20}, the \emph{Carter strategy} $\C_h$ on $h$ hats is
obtained from $\C_1 = \str{0}$ by double shifts for odd $h$, and by one
more black shift for even $h$; so $\C_3 = \str{01002120}$, and in terms
of substitutions $\C_{h+2} = \C_3[\C_h]$ and
$\C_{h+1} = \str{0100}[\C_h]$ for odd $h$. The name is that of the posts
of 2020, not of the literature. Strategies with these scores go back to 2014: Wagon's notes \cite{Wa14}
give them on up to eight hats, with the recursion $x \mapsto 16x+6$, and
Theorem~76.2 of \cite{VW20} states them for every $h$, with a sketch of
proof, crediting Carter, Reyes, Rosenberg and Tiefenbruck for odd $h$ and
Wagon and Zimmermann for even $h$; \cite{BCCZW} builds them again for odd
$h$. Our construction for even $h$ (one more shift) differs from that of
\cite{VW20} (a reset that overlaps by one hat).

\begin{corollary}\label{cor:carter}
For every $h \ge 1$, $|H(\C_h)| = \lfloor 7\cdot 4^h/20 \rfloor$: the
win rate of $\C_h$ is $\frac{7}{20} - \frac{2}{5\cdot 4^h}$ for odd $h$
and $\frac{7}{20} - \frac{3}{5\cdot 4^h}$ for even $h$.
\end{corollary}

(Lean: \lean{hits\_carter}; proof in Appendix~\ref{app:carter}.) A win rate on $h$ hats is a multiple of
$4^{-h}$, and $7\cdot 4^h/20$ is never an integer, so
$\lfloor 7\cdot 4^h/20\rfloor$ is the largest number of hits that does
not pass $7/20$. This gives the rate at which finite stacks approach
$7/20$, improving the bound $\frac{7}{20} - C_\varepsilon
(4^{1-\varepsilon})^{-h}$ of \cite{BCCZW}; with their remark that
$|\frac{7}{20} - V_{2,h}| \ge \frac{1}{5\cdot 4^h}$ for the best win
rate $V_{2,h}$ on $h$ hats, the order $4^{-h}$ cannot be improved if the
conjecture holds.

The strategy $\C_h$ is also given by a rule that does not depend on $h$
(Lean: \lean{carter\_words}). If the bottom hat seen is black, point at
the lowest white hat when it is at an odd position, and at the bottom hat
otherwise. If the bottom hat seen is white, look at the lowest black hat:
at an even position, point at it; at an odd position, point at the hat
just above it if that hat exists and is black; in every other case point
at the bottom hat.

\begin{corollary}\label{cor:conj}
The following are equivalent: (i) no strategy on any number of hats has a
win rate above $7/20$; (ii) for every $h \ge 1$ the best number of hits
on $h$ hats is $\lfloor 7\cdot 4^h/20\rfloor$.
\end{corollary}

(Lean: \lean{conjecture\_iff}.) That (i) implies (ii) is a remark of
Yeroshkin reported in \cite{VW20}. Statement (i) is the conjecture
$V_2 = 7/20$ for plans in which both players use the same strategy.

\begin{corollary}\label{cor:gain}
An operation $m$ that adds $k$ hats has a gain of at most $B_{k+1} - 4^k$,
where $B_j$ is the best number of hits on $j$ hats. With $B_2 = 5$,
$B_3 = 22$, $B_4 = 89$: at most $1$, $6$ and $25$ for one, two and three
hats added.
\end{corollary}

(Lean: \lean{hits\_sub\_le}; $B_2$ and $B_3$ are proved by the kernel,
$B_4 = 89$ by a certified search run as a compiled program; $B_5 = 358$
is the answer of solvers and of the enumeration of
Section~\ref{sec:programs}.) In particular no operation built from
resets, shifts or any substitution can do better than the Carter
strategies unless some strategy on few hats already does.

\subsection{A second operation that adds two hats}

The double shift is $m[s]$ for $m = \C_3$. The best strategies on three
hats are all equivalent to $\C_3$ (Section~\ref{sec:programs}), but
equivalent strategies $m$ do not give equivalent operations, since the
bottom hat of $m$ plays a special part. The \emph{twisted double shift}
is $m[s]$ for $m = \str{02110200}$. In words: play $s$ on the $h$ lowest
hats, except that $h$ black hats with a white hat just above send to the
top hat, and $h$ hats of a single color with a black hat just above and
a white hat at the top send to the hat just above.

\begin{proposition}\label{prop:twisted}
The twisted double shift of $s$ has $16|H(s)|+6$ hits. For each of the
sixteen strategies $s$ on two hats, it is not equivalent to the double
shift of $s$, whatever the values on the all-white and all-black stacks.
\end{proposition}

(Lean: \lean{hits\_tds}, \lean{tds\_eq\_sub},
\lean{tds\_not\_renamed\_ds}; proof in Appendix~\ref{app:twisted}.) On five hats, up to equivalence, the
exchange of the colors and two stacks on which the choice is free, the
best strategies are exactly two: $\C_5$, and the twisted double shift of
$\C_3$ (Section~\ref{sec:programs}).

\section{Biased coins}\label{sec:bias}

Let now each hat be black with probability $p$, independently, and
$q = 1-p$. The \emph{value} $\mathrm{val}_p(s)$ of a strategy is the
probability that both players point at a black hat, and $V_2(p)$ the
supremum of the values. In \cite{BFG} four strategies $S_1,\ldots,S_4$
are given, all of value $7/20$ at $p = 1/2$; \cite{BCCZW} adds a fifth,
$S_5$, better than these for $p < 0.312$.

\begin{proposition}[Exchange of the colors]\label{prop:dual}
For every strategy $s$ and every $p$,
$\mathrm{val}_{1-p}(\bar s) = \mathrm{val}_p(s) + 1 - 2p$. Hence
$V_2(1-p) = V_2(p) + 1 - 2p$.
\end{proposition}

\begin{proof}
A player points at a black hat with probability $p$ whatever the
strategy, since the choice does not depend on the player's own hats. So
both point at a white hat with probability $1 - 2p + \mathrm{val}_p(s)$,
and exchanging the colors turns this event into a win.
\end{proof}

(Lean: \lean{val\_dual}.) In \cite{BFG} the strategies $S_1$ and $S_3$,
which give the bounds for $p \le 1/2$ and $p \ge 1/2$, are stated
separately; they are dual to each other.

The substitution has a formula for every $p$. For a strategy $m$ on
$k+1$ hats let $A_{cc'}$ be the probability, over the $k$ upper hats of
both players, that the pair of stacks with bottom hats of colors $c$ and
$c'$ is a hit of $m$, and $a$ the probability that $m$ points at the
bottom hat when the bottom hat seen is black.

\begin{theorem}\label{thm:valsub}
For $s$ on $h \ge 1$ hats and $m$ on $k+1$ hats,
\begin{align*}
  \mathrm{val}_p(m[s]) = {} & \mathrm{val}_p(s) - p^{2h}
    + q^{2h} A_{00} + (pq)^h (A_{01} + A_{10}) + p^{2h} A_{11} \\
  & {} - 2p^h (1-a)\, pq\, \bigl(q^{h-1} - p^{h-1}\bigr) .
\end{align*}
\end{theorem}

(Lean: \lean{val\_sub}; proof in Appendix~\ref{app:valsub}.) The last term vanishes at $p = 1/2$, where the
formula is Theorem~\ref{thm:sub}. For $p < 1/2$ it is a loss unless
$a = 1$: a player should not leave the lowest hats on seeing them all
black. The resets of \cite{BFG} and \cite{BCCZW} leave on both colors.

\begin{theorem}[White resets]\label{thm:white}
For $0 < p \le 1$,
\[
  V_2(p) \;\ge\; \frac{3p^2 - 6p^3 + 8p^4 - 6p^5 + 2p^6}{1 - (1-p)^6} .
\]
\end{theorem}

\begin{proof}
A white reset adds values: the white reset of $s$ (on $h$ hats) by $b$
has the value $\mathrm{val}_p(s) + q^{2h}\,\mathrm{val}_p(b)$, since the
pairs in which exactly one player has $h$ white hats below are lost with
or without the reset. The numerator is $\mathrm{val}_p(\C_3)$. The
white reset of $\C_3$ by itself $j$ times has the value
$\mathrm{val}_p(\C_3)\,(1 + q^6 + \cdots + q^{6j})$.
\end{proof}

(Lean: \lean{val\_lr}, \lean{val\_c3}, \lean{val\_whiteChain},
\lean{tendsto\_whiteChain}.) In words: play $\C_3$ on the first block of
three hats that is not all white. For fair coins this strategy is in
Wagon's notes \cite{Wa14}, where Roche gives it with the value $22/63$;
what is new here is its value for $p \ne 1/2$ and the comparison that
follows. At $p = 1/2$ the bound is $22/63$,
below $7/20$. Table~\ref{tab:bias} compares it with the best of
$S_1,\ldots,S_5$; it is larger at every $p = i/20$ for $1 \le i \le 8$
and, by Proposition~\ref{prop:dual} and the mirror strategy, for
$12 \le i \le 19$. (The comparison is by exact rational arithmetic in a
program, \texttt{biased.py} in the repository.)

\begin{table}[ht]
\centering
\begin{tabular}{@{}ccc@{}}
\toprule
$p$ & best of $S_1,\ldots,S_5$ & white resets of $\C_3$ \\
\midrule
$1/10$ & $0.052755$ & $24742/468559 = 0.052804$ \\
$1/5$  & $0.112240$ & $1297/11529 = 0.112499$ \\
$3/10$ & $0.180491$ & $17742/98039 = 0.180969$ \\
$2/5$  & $0.259458$ & $967/3724 = 0.259667$ \\
$9/20$ & $0.303257$ & $0.303080$ \\
$1/2$  & $7/20$ & $22/63$ \\
\bottomrule
\end{tabular}
\caption{Lower bounds on $V_2(p)$.}\label{tab:bias}
\end{table}

With blocks of four and five hats in place of three, the white resets do
less well for the values of $p$ that we tried, and at $p = 3/10$ the
best strategy that a search finds on six hats is within $3\cdot 10^{-6}$
of two blocks of white resets.

\section{More players}\label{sec:players}

Let there be $n \ge 2$ players. A \emph{plan} on $T$ hats gives each
player a position of their own stack as a function of the stacks of
everybody; it is \emph{blind} when no player uses their own stack, which
is the rule of the game. Its \emph{wins} are the placements of the hats,
out of $2^{nT}$, for which all the hats pointed at are black. Players may
have different tables; the tables of \cite{BFG} are symmetric, one for
all.

\begin{definition}
The \emph{reset} of a plan $S$ on $T$ hats by a plan $M$ on $k$ hats: a
player who sees a stack that is all white on its $T$ lowest hats, or
sees only stacks that are all black there, plays $M$ on the $k$ hats
above; any other player plays $S$ on the $T$ lowest hats.
\end{definition}

This is the reset of \cite{BFG}, where the limit below is stated with a
proof in words.

\begin{theorem}\label{thm:reset}
Let $S$ and $M$ be blind plans on $T \ge 1$ and $k \ge 1$ hats, with
$W_S$ and $W_M$ wins, and $P = 2^T$. The reset of $S$ by $M$ has
\[
  2^{nk}\,(W_S - 1) + R\,W_M
  \quad\text{wins, where}\quad
  R = P^n - (P-1)^n - n(P-1)^{n-1} + n + 1 .
\]
If $T \ge 2$, the win rates of $S$ reset by itself again and again tend
to
\[
  \frac{W_S - 1}{(P-1)^n + n(P-1)^{n-1} - (n+1)} .
\]
\end{theorem}

\begin{proof}[Sketch of proof]
Sort the ways to choose the $T$ lowest hats of everybody. (a) Everybody
goes up: this happens when two players at least are all white, or one is
all white and the others all black, or all are all black; there are $R$
such ways, and each is worth $W_M$. (b) One player is mixed and all the
others are all black: the others stay and point at black hats; the mixed
player goes up alone and points at a black hat for half of the hats
above, as often as by staying, where the hat pointed at is black for
$2^{T-1}-1$ of the $2^T - 2$ mixed stacks. (c) All the other ways:
either a player who is all white stays, and the game is lost with $S$
too, or nobody goes up and $S$ decides. So every way outside (a) is
worth $2^{nk}$ times what it is worth for $S$, and $S$ wins exactly one
of the ways of (a), the one where all are all black. The limit follows
from the recursion on the win rates, whose ratio $R/P^n$ is less than
$1$ for $T \ge 2$.
\end{proof}

(Lean: \lean{winsN\_reset}, \lean{resetCard\_add}, \lean{tendsto\_rateN},
\lean{limitN\_eq}; complete proof in Appendix~\ref{app:reset}.) For $n = 2$, $R = 4$ and the limit is
$(W_S-1)/(4^T-4)$: the reset is the double reset of
Section~\ref{sec:two}, and the limit is $7/20$ for $S = \C_T$ with $T$
odd (Lean: \lean{tendsto\_rate\_drChain\_carterOdd}). The substitution
of Theorem~\ref{thm:sub} also makes sense for three players, with
$8^k(W_S-1)+W_M$ wins for $M$ on $k+1$ hats (Lean: \lean{wins3\_sub});
for three players the reset does better.

\subsection{New bounds}

Theorem~\ref{thm:reset} says where to look. For three players, the best
table known on five hats has $9120$ wins out of $32768$ \cite{BFG}; we
found one with the same score. Reset by one hat, it gives a table on six
hats with $73050$ wins, and a table on six hats improves the bound of
\cite{BFG} as soon as it has $73117$ wins. A simulated annealing on
symmetric tables that starts from the first of these, and changes one
entry at a time, passes that number in a few minutes, whereas the same
search from a random table stopped at $71819$ in our runs. We repeated this from
six to seven hats and from seven to eight; for four players, random
starts on four hats suffice.

\begin{theorem}\label{thm:bounds}
$V_3 \ge 73225/261950 = 0.279538\ldots$ and
$V_4 \ge 15191/64120 = 0.236915\ldots$
\end{theorem}

(Lean: \lean{tendsto\_six}, \lean{tendsto\_four4}.) The proof is a
symmetric table on six hats with $73226$ wins for three players, a
symmetric table on four hats with $15192$ wins for four players, and
Theorem~\ref{thm:reset}. Both tables are in the Lean sources, and the
kernel counts their wins. Table~\ref{tab:players} lists all the tables;
those on seven and eight hats, too large for the kernel as we encoded
them, are counted by two programs written independently.

\begin{table}[ht]
\centering
\begin{tabular}{@{}cccccl@{}}
\toprule
players & hats & wins & out of & limit & counted by \\
\midrule
3 & 5 & 9120 & $2^{15}$ & 0.279124 & \cite{BFG} \\
3 & 6 & 73226 & $2^{18}$ & 0.279538 & Lean \\
3 & 7 & 586263 & $2^{21}$ & 0.279603 & two programs \\
3 & 8 & 4690902 & $2^{24}$ & 0.279612 & two programs \\
4 & 4 & 14845 & $2^{16}$ & 0.231503 & \cite{BFG} \\
4 & 3 & 886 & $2^{12}$ & 0.234873 & Lean \\
4 & 4 & 15192 & $2^{16}$ & 0.236915 & Lean \\
4 & 5 & 246459 & $2^{20}$ & 0.236370 & two programs \\
\bottomrule
\end{tabular}
\caption{Symmetric tables and the limit of their resets.}
\label{tab:players}
\end{table}

The notes \cite{Wa14} of 2014 mention a table on four hats with $1125$
wins for three players and say that it ``augments to $64/225$ with the
reset strategy'', which would be $0.2844$. By Theorem~\ref{thm:reset} a
table on four hats with $1125$ wins gives $1124/4046 = 0.2778$, and
\cite{BFG}, written later by the same group, gives $0.279124$ as the best
bound known; we take the $64/225$ for a slip.

Two remarks. A table on three hats already beats the bound of \cite{BFG}
for four players: resets weigh more on few hats. And for three players
each hat added brings a little more; nothing indicates that these bounds
are close to $V_3$ and $V_4$. For two players the same search, started
from $\C_7$ and $\C_8$, finds nothing better than the Carter strategies.

\section{Remarks on upper bounds}\label{sec:upper}

The upper bounds on $V_2$ of \cite{BFG} and \cite{Ya26} come from a
\emph{hint}: Alice's stack is known to be one of the $2(2^d-1)$
non-constant affine functions on $\mathbb{F}_2^d$, her hats being its
values at the $2^d$ points. Write $H_d$ for the value of that game;
$V_2 \le H_d$, with $H_3 = 81/224$ \cite{BFG} and $H_4 = 57/160$
(\cite{Ya26}; we confirmed this value and its $11$ classes of optimal
solutions with a program written independently).

\begin{proposition}\label{prop:hint}
For any two strategies $a$, $b$ on $d+1$ hats, used by Alice and Bob,
the number of hits is at most $2^d + 2^{d+1}(2^{d+1}-2)\,H_d$.
\end{proposition}

(Lean: \lean{hitsTwo\_le\_hintWins}; proof in Appendix~\ref{app:hint}.) The values of an affine function at
the points $0, e_1, \ldots, e_d$ run over the stacks of $d+1$ hats that
are not of a single color, once each. With $\C_{d+1}$ for both players
the inequality is an equality for $d = 1,\ldots,4$, and the best values
found by annealing in \cite{Ya26} for $d = 5, 6$ are again those of
$\C_6$ and $\C_7$. So $H_d \ge \frac{7}{20} + \Theta(2^{-d})$; if equality persists for
every $d$, then $H_d \to 7/20$ and the conjecture follows. With
$H_4 = 57/160$ the proposition gives $358$ as the best score on five
hats, also for two different strategies, which answers a question of
\cite{Wa14}; the same question for $1433$ on six hats is open.

Not every optimal solution of the hint game is a strategy of the real
game: of the $11$ classes for $d = 4$, four need Alice to point at
points that are affinely dependent. The hint game can be read as
follows. Bob's hats decide which vectors $v \ne 0$ are active and a
target bit $T(v)$ for each; Alice looks for a point $x$ with
$v\cdot x = T(v)$ for as many active $v$ as she can, and the value is
$\frac{3}{8}$ less the expected number of equations she gives up,
divided by $2(2^d-1)$. Three active vectors $v$, $w$, $v+w$ whose
targets do not add up force her to give up one. For all the optimal
solutions with $d = 4$, and for the solutions that come from $\C_{d+1}$
with $d \le 6$, these triples are the only obstruction (a result of
programs). The conjecture would follow from a proof that the expected
number of equations given up is at least $2^d/20 - o(2^d)$.

\section{The best strategies on few hats}\label{sec:programs}

The statements of this section are results of programs. An exhaustive
search gives the best scores $1$, $5$, $22$, $89$ on one to four hats
(the last two are also certified in Lean), and $358$ on five hats by
three programs and by Proposition~\ref{prop:hint}. Up to equivalence and
to the values on the all-white and all-black stacks, the best strategies
form $2$, $1$, $10$ and $27$ classes on two to five hats. On four hats,
six classes are substitutions of smaller best strategies and four are
not; these four are one strategy up to the exchange of the colors and two
free choices. On five hats the $27$ classes are the two strategies of
Section~\ref{sec:two}, each with two stacks on which three choices are
equally good. On six hats, $2392$ runs of simulated annealing found
$1971$ different strategies with $1433 = |H(\C_6)|$ hits and none with
more; each of them is, after free choices, the substitution of a best
strategy on three hats into a best strategy on four hats.

\section{Formal verification, data and code}\label{sec:lean}

The Lean~4 sources (directory \texttt{formal/hats} of \cite{web}, built
against Mathlib) contain every statement cited with a Lean name. An
audit file prints the axioms of each theorem; all of them use the three
standard axioms of Lean only, except the best score on four hats, which
is obtained by running a certified search as a compiled program. For
Theorem~\ref{thm:bounds} the tables are literals of the Lean sources and
the counts are done by \texttt{decide} in the kernel.

\begin{table}[ht]
\centering
\small
\begin{tabular}{@{}ll@{}}
\toprule
statement & Lean \\
\midrule
Theorem~\ref{thm:sub} & \lean{hits\_sub} \\
Proposition~\ref{prop:instances} & \lean{lr\_eq\_sub}, \lean{rr\_eq\_sub},
  \lean{dr\_eq\_sub}, \lean{ls\_eq\_sub}, \lean{rs\_eq\_sub},
  \lean{ds\_eq\_sub} \\
Corollary~\ref{cor:carter} & \lean{hits\_carter}, \lean{carter\_words} \\
Corollary~\ref{cor:conj} & \lean{conjecture\_iff} \\
Corollary~\ref{cor:gain} & \lean{hits\_sub\_le} \\
Proposition~\ref{prop:twisted} & \lean{hits\_tds},
  \lean{tds\_not\_renamed\_ds} \\
Proposition~\ref{prop:dual} & \lean{val\_dual} \\
Theorem~\ref{thm:valsub} & \lean{val\_sub} \\
Theorem~\ref{thm:white} & \lean{tendsto\_whiteChain}, \lean{val\_c3} \\
Theorem~\ref{thm:reset} & \lean{winsN\_reset}, \lean{resetCard\_add},
  \lean{tendsto\_rateN}, \lean{limitN\_eq} \\
Theorem~\ref{thm:bounds} & \lean{tendsto\_six}, \lean{tendsto\_four4} \\
Proposition~\ref{prop:hint} & \lean{hitsTwo\_le\_hintWins} \\
\bottomrule
\end{tabular}
\caption{The statements and their Lean names.}
\end{table}

The tables of Table~\ref{tab:players} are text files of the repository
(one digit per entry). The first sixteen hexadecimal digits of their
SHA-256 sums are: six hats, \texttt{a2929fc0f57e9cbc}; seven hats,
\texttt{567dd9eeed59d44b}; eight hats, \texttt{328f8a1d8aab5d51}; four
players on three hats, \texttt{32a068d81a5a9c99}; four players on four
hats, \texttt{3ab3840ccd415040}. The search programs, the plain scorers
used for the recounts, and a record of the literature consulted are in
the same repository.

\section{Open questions}

\begin{question}
For small $p$, is $V_2(p)$ equal to the bound of
Theorem~\ref{thm:white}?
\end{question}

\begin{question}
For three players and more, is there a rule to leave the lowest hats
that does better than the reset of Theorem~\ref{thm:reset}, with tables
made for it? Do the bounds of Table~\ref{tab:players} keep growing with
the number of hats, and toward what?
\end{question}

\begin{question}
On every number of hats, is every best strategy for two players a
substitution of best strategies on fewer hats, up to equivalence, colors
and free choices, with finitely many exceptions such as the one on four
hats?
\end{question}

\begin{question}
Is $H_d$ the value given by $\C_{d+1}$ in Proposition~\ref{prop:hint}
for every $d$? Is the obstruction by triples of Section~\ref{sec:upper}
the only one at the optimum?
\end{question}


\appendix
\section{Complete proofs}\label{app:proofs}

\subsection{Proposition~\ref{prop:basic}}

(1) A pair $(t,u)$ with $t$ all white is never a hit, since
$t_{s(u)} = 0$. The value of $s$ on the all-white stack is used only for
such pairs. Let $b$ be the all-black stack. The value $s(b)$ is used for
the pairs $(t,b)$ and $(b,t)$, which are hits exactly when
$t_{s(b)} = 1$: this holds for $2^{h-1}$ stacks $t$ whatever $s(b)$ is.

(2) A renaming $\pi$ of the positions acts on stacks and on strategies,
by $s^\pi(\pi t) = \pi(s(t))$, and $(t,u) \mapsto (\pi t, \pi u)$ is a
bijection from the hits of $s$ to those of $s^\pi$.

(3) For a fixed $u$, the stacks $t$ with $t_{s(u)} = 1$ are half of all
stacks. So the pairs in which Alice points at a black hat are $4^h/2$,
and so are those in which Bob does. The pairs in which both point at a
white hat are therefore $4^h - 2\cdot 4^h/2 + |H(s)| = |H(s)|$, and
exchanging the colors turns them into the hits of $\bar s$.

\subsection{Corollary~\ref{cor:carter}}\label{app:carter}

The strategies $\C_3 = \str{01002120}$ and $\str{0100}$ have $22$ and
$5$ hits, by a direct count. By Theorem~\ref{thm:sub},
$|H(\C_{h+2})| = 16(|H(\C_h)|-1) + 22 = 16|H(\C_h)| + 6$ and
$|H(\C_{h+1})| = 4(|H(\C_h)|-1)+5 = 4|H(\C_h)|+1$ for odd $h$. With
$|H(\C_1)| = 1$, induction gives $20|H(\C_h)| = 7\cdot 4^h - 8$ for odd
$h$, and then $20|H(\C_h)| = 7\cdot 4^h - 12$ for even $h$. Since
$7\cdot 4^h$ is congruent to $8$ modulo $20$ for odd $h$ and to $12$ for
even $h \ge 2$, these are the integer parts of $7\cdot 4^h/20$.

\subsection{Proposition~\ref{prop:twisted}}\label{app:twisted}

The strategy $m = \str{02110200}$ on three hats has $22$ hits, so by
Theorem~\ref{thm:sub} the operation $s \mapsto m[s]$ has the gain
$22 - 16 = 6$. Reading Definition~\ref{def:sub} for this $m$: the stacks
of three hats with a black bottom hat are $1,3,5,7$, with
$m = 2,1,2,0$, and those with a white bottom hat are $0,2,4,6$, with
$m = 0,1,0,0$; this is the description in words. For the second claim
there are sixteen strategies on two hats and twenty-four renamings of
four positions; for each of them some stack that is not of a single
color tells the two strategies apart. This finite check is done in Lean.

\subsection{Theorem~\ref{thm:valsub}}\label{app:valsub}

Write $w(t)$ for the probability of a stack $t$, so that
$\mathrm{val}_p(s) = \sum w(t)w(u)$ over the hits $(t,u)$ of $s$. The
all-white stack of $h$ hats has probability $q^h$ and the all-black one
$p^h$; the mixed stacks have total probability $1 - q^h - p^h$, and
among them those with a black hat at a given position have total
probability $p - p^h$.

Sort the pairs as in the proof of Theorem~\ref{thm:sub}. \emph{Both
mixed:} the hats above have total probability $1$, so these pairs
contribute to $m[s]$ what they contribute to $s$, say $M$. \emph{Both of
a single color}, $c$ and $c'$: the lower parts have probability
$w_c w_{c'}$ with $w_0 = q^h$, $w_1 = p^h$, and the pair is a hit when
the stacks of $k+1$ hats are a hit of $m$; this gives
$q^{2h}A_{00} + (pq)^h(A_{01}+A_{10}) + p^{2h}A_{11}$. \emph{All white
and mixed:} nothing. \emph{All black and mixed:} the player who sees the
mixed stack points at a black hat. The other one, seeing the all-black
lower part under hats $x$, either stays (probability $a$ over $x$) and
points at a fixed lower position of a mixed stack, which is black with
total probability $p - p^h$, or goes up (probability $1-a$) to a
position above, which is black with probability $p$, the mixed lower
part having total probability $1-q^h-p^h$. With the two orders:
$2p^h\bigl(a(p-p^h) + (1-a)\,p\,(1-q^h-p^h)\bigr)$.

For $s$ alone the same sorting gives
$\mathrm{val}_p(s) = M + 2p^h(p - p^h) + p^{2h}$. Subtracting, and
using $(p-p^h) - p(1-q^h-p^h) = pq(q^{h-1}-p^{h-1})$, gives the formula.

For Theorem~\ref{thm:white}, the white reset of $s$ by $b$ is $m[s]$ with
$m$ the white reset of $\str{0}$ by $b$: then $a = 1$,
$A_{00} = \mathrm{val}_p(b)$, $A_{01} = A_{10} = 0$ and $A_{11} = 1$, so
that the gain is $q^{2h}\,\mathrm{val}_p(b)$. The value of
$\C_3$ is a sum over its $22$ hits:
$3p^2q^4+6p^3q^3+8p^4q^2+4p^5q+p^6$, which is the numerator of the
theorem.

\subsection{Theorem~\ref{thm:reset}}\label{app:reset}

Call a stack \emph{white}, \emph{black} or \emph{mixed} according to its
$T$ lowest hats. A player goes up when some other player is white, or
all the others are black.

\emph{Who goes up.} If two players at least are white, everybody sees a
white player and goes up. If exactly one player is white, the others go
up, and the white player goes up exactly when all the others are black.
If nobody is white, a player goes up exactly when all the others are
black; if all go up, take two players: all but the first are black and
all but the second are black, so all are black (here $n \ge 2$ is
used). Hence everybody goes up in exactly three cases, (a1) at least two
white players, (a2) one white player and the others black, (a3) all
black; the number of ways to choose the lowest hats in these cases is
$\bigl(P^n - (P-1)^n - n(P-1)^{n-1}\bigr) + n + 1 = R$.

\emph{What each way is worth.} In the cases (a) all the players play $M$
on the hats above, for $W_M$ wins. Outside (a) there are three kinds.
(b) Nobody is white, one player is mixed and the others are black: the
mixed player goes up alone. The others stay, and point at black hats
since their lowest hats are all black. The mixed player points at a
position above that depends on the hats above of the others only, so at
a black hat for half of the $2^{nk}$ ways to choose the hats above.
(c1) Exactly one player is white and the others are not all black: the
white player stays, and points at a white hat. (c2) Nobody is white and
two players at least are not black: nobody goes up, and the result is
that of $S$ on the lowest hats, for each of the $2^{nk}$ ways above.

\emph{The same ways, for $S$.} In (a1) and (a2) $S$ loses, a white
player being present; in (a3) it wins. In (b), with the mixed player $i$
and its lowest hats $\ell$, $S$ wins when $\ell$ is black at the
position that $S$ gives to player $i$, which does not depend on $\ell$:
$2^{T-1}-1$ of the $P-2$ mixed stacks, that is half. In (c1) $S$ loses.
So over the ways of (b) the reset has $2^{nk}$ times the wins of $S$,
and in (c1), (c2) each way has $2^{nk}$ times its wins for $S$. Summing,
the reset has $R\,W_M + 2^{nk}(W_S - 1)$ wins.

\emph{The limit.} Let $r_j$ be the win rate of $S$ reset by itself $j$
times, on $T(j+1)$ hats. The count gives
$r_{j+1} = (W_S-1)/P^n + (R/P^n)\,r_j$. For $T \ge 2$ one has
$(P-1)^{n-1} \ge 3$, so $R < P^n$, and $r_j$ tends to the fixed point
$(W_S-1)/(P^n - R)$.

\subsection{Proposition~\ref{prop:hint}}\label{app:hint}

In the hint game of dimension $d$, Alice's stack is $f$, uniform among
the $2(2^d-1)$ non-constant affine functions on $\mathbb{F}_2^d$, her
hats being indexed by the points; Bob has fair hats, as many as he
wants; Bob sees $f$ and Alice sees Bob's hats; they win when both point
at a black hat. $H_d$ is the best probability of winning.

Let $a$, $b$ be strategies on $d+1$ hats for Alice and Bob. The map
\[ f \longmapsto t_f = \bigl(f(0), f(e_1), \ldots, f(e_d)\bigr) \]
is a bijection from the non-constant affine functions onto the stacks of
$d+1$ hats that are not of a single color. In the hint game let Bob, who sees $f$, point at
$b(t_f)$ among his $d+1$ first hats, and let Alice, who sees Bob's stack
$u$, point at the point $0$ or $e_i$ according to $a(u)$. They win
exactly when $(t_f,u)$ is a hit of the pair $(a,b)$. So the hits of
$(a,b)$ in which Alice's stack is not of a single color are at most
$2^{d+1}(2^{d+1}-2)H_d$. If Alice's stack is all white there is no hit;
if it is all black, the pair is a hit when Bob's stack is black at the
position that Bob chooses on seeing the all-black stack: $2^d$ stacks.

\section*{Declaration on the use of AI tools}

The operations (resets and shifts), their formulas and the Carter
strategy for every number of hats are from blog posts that the author
wrote alone in 2020 \cite{R20}. The present work was done in November 2026 by the
author in collaboration with Claude, an AI system made by Anthropic. The
collaboration covered the mathematics (the substitution theorem and its
consequences, the bounds for three and four players and for biased
coins, the remarks on upper bounds), the programs, the formal proofs and
the writing of this text, most of which was drafted by Claude. The
author set the questions, directed the work, read and questioned the
results, and is the sole author of the paper: the author takes
responsibility for its content. The statements checked in Lean do not
rest on trust in either the author or the assistant; the correspondence
between the Lean statements and the text, and the results of programs,
rest on the usual reading.

\emergencystretch=3em
\begin{thebibliography}{99}

\bibitem{AFKK} N.~Alon, E.~Friedgut, G.~Kalai and G.~Kindler, The
success probability in Levine's hat problem, and independent sets in
graphs, \emph{SIAM J. Discrete Math.} \textbf{37} (2023),
arXiv:2208.06858.

\bibitem{BCCZW} C.~Bouquet, S.~Chikhi, T.~Charles, Y.~Zhou and E.~Wang,
An analytical framework for the Levine hats problem: new strategies,
bounds and generalizations, preprint, 2025, arXiv:2508.01737.

\bibitem{BFG} J.~Buhler, C.~Freiling, R.~Graham, J.~Kariv, J.~R.~Roche,
M.~Tiefenbruck, C.~Van~Alten and D.~Yeroshkin, On Levine's notorious hat
puzzle, \emph{Integers} \textbf{21A} (2021), \#A8, arXiv:1407.4711.

\bibitem{FKK21} E.~Friedgut, G.~Kalai and G.~Kindler, The success
probability in Lionel Levine's hat problem is strictly decreasing with
the number of players, preprint, 2021, arXiv:2103.01541.

\bibitem{HT25} S.~Heilman and O.~Tamuz, A Fourier approach to Levine's
hat puzzle, preprint, 2025, arXiv:2503.09042.

\bibitem{Kh11} T.~Khovanova, How many hats can fit on your head?, blog
post, 2011,
\url{https://blog.tanyakhovanova.com/2011/04/how-many-hats-can-fit-on-your-head/}.

\bibitem{R20} F.~Ruget, Towers of hats (1/3), Hats calculator (2/3),
Infinitely many hats (3/3), three blog posts, 30 August, 5 September and
26 September 2020, first published at \texttt{blog.atlant.is}. The
Internet Archive holds the text of the three posts as captured on
28 September 2020:
\url{https://web.archive.org/web/20200928125658/https://blog.atlant.is/?cat=16}.
They are kept unchanged at
\url{https://douzebis.github.io/blog.atlant.is/versions/the-game-of-hats/2020/2020-08-30-tower-of-hats/}
and in the directory \texttt{content/versions/the-game-of-hats/2020} of
\repo.

\bibitem{VW20} D.~J.~Velleman and S.~Wagon, \emph{Bicycle or Unicycle? A
Collection of Intriguing Mathematical Puzzles}, MAA Press, 2020,
Problem~76.

\bibitem{Wa14} S.~Wagon, Problem 1179: An infinitely puzzling hat
problem, and its solution notes, Macalester Problem of the Week, 2014,
\url{https://stanwagon.com/potw/current_solutions/s1179.html}.

\bibitem{web} F.~Ruget, The game of hats, web article, 2026 (a revision of
\cite{R20}), with the Lean sources, programs and tables of the present
paper, \repo.

\bibitem{Ya26} Yang Song, Exact referee relaxations for Levine's
two-player hat problem, note with programs, 2026,
\url{https://github.com/yangsong97/levine-hat-bounds}.

\end{thebibliography}

\end{document}
